% Section 5: The natural index and the expansion to n^{-6} (Theorem A = Theorem 5.7(iii)).
\section{The natural index and the expansion of the lattice}\label{sec:natural}

This section proves Theorem~\ref{thm:A}. The $1/n$ term and the $(s-1)^2$ part of the $n^{-2}$ term of the lattice are
one index shift (\S\ref{sec:shift}); for Bessel zeros the shift is exact to all orders (\S\ref{sec:bessel}). We then
extend the analysis of \cite{PaperII} and Theorem~\ref{thm:E2E3} to every order $K$ in $1/n$ and every order $B$ in the
band-factor parameters (\S\S\ref{sec:recursion}--\ref{sec:N6}). The analytic part is the order-$K$ version of
estimates already proved at order $2$; the algebraic part is a finite residue recursion, which we evaluate in exact
arithmetic. Strategy: the band factor is a tilt (Lemma~\ref{lem:Fprime}), tilts act on the jump data by an entire gauge
(Lemma~\ref{lem:gauge}), so every coefficient of the expansion is an entire function of the tilt parameters $\mu$ with
$n$-independent Taylor coefficients, and the Taylor coefficients are computed by the recursion.

\subsection{The index shift}\label{sec:shift}
Put $\delta_N:=\frac{\alpha+1}2$, so that $N=n+\delta_N$, and
\begin{equation}\label{eq:f2}
  f_2=f_2(p,\alpha):=\frac{p+1}{12}-\frac{(p+2)\alpha^2}8 .
\end{equation}

\begin{proposition}\label{prop:N1}
For all $p,s$, $d_1=-p\delta_N$ and $d_2(p,s)=\frac{p(p+1)}2\delta_N^2+f_2$. Equivalently,
$1+d_1/n+d_2/n^2=(n/N)^p(1+f_2N^{-2})+O(n^{-3})$.
\end{proposition}

\begin{proof}
With $s-1=p(\alpha+1)$ and $s-1-p=p\alpha$, the three terms of $d_2$ in Theorem~\ref{thm:E2E3} are
$p(p+1)(\alpha+1)^2/8=\frac{p(p+1)}2\delta_N^2$, $-(p+2)\alpha^2/8$ and $(p+1)/12$. Since
$(1+\delta/n)^{-p}=1-p\delta/n+\frac{p(p+1)}2\delta^2n^{-2}+O(n^{-3})$, the claim follows.
\end{proof}

So the lattice expansion to order $n^{-2}$ is an index shift by $\delta_N$ followed by a bare function with no $1/N$
term. The shift has two parts: $\frac\alpha2$, attached to the hard edge of Bessel order $\alpha$, and
$\frac12=-\zeta(0)$, the regularised count of the nodes (Section~\ref{sec:conj}).

\subsection{Bessel zeros}\label{sec:bessel}
For $\mu>-1$ let $j_{\mu,1}<j_{\mu,2}<\dotsb$ be the positive zeros of the Bessel function $J_\mu$ and
$\lambda_k:=(j_{\mu,k}/\pi)^2$. For $p=s=2$ the measure~\eqref{eq:sigmaL} is
$\sigma_\Lambda=\pi^2\sum_kj_{\mu,k}^{-2}\delta_{\pi^2/j_{\mu,k}^2}$, and $\mu=\frac12$ gives the lattice $\lambda_k=k^2$.

\begin{proposition}\label{prop:bessel}
The recurrence coefficients of $\sigma_\Lambda$ satisfy, for $n\ge1$, with $N:=n+\mu/2$,
\[
  \frac{n^2b_n}{c(2)}=\frac{4n^2}{(2n+\mu)\sqrt{(2n+\mu)^2-1}}=\Bigl(\frac nN\Bigr)^2\Bigl(1-\frac1{4N^2}\Bigr)^{-1/2}.
\]
The bare function $(1-\frac1{4N^2})^{-1/2}=1+\frac1{8N^2}+\frac3{128N^4}+\frac5{1024N^6}+\dotsb$ is even in $1/N$, and its
coefficients are $f_2(2,-\frac12)$, $f_4(2,-\frac12)$, $f_6(2,-\frac12)$ of~\eqref{eq:f2}, \eqref{eq:f4}, \eqref{eq:f6}.
\end{proposition}

\begin{proof}
The Lommel polynomials of order $\nu=\mu+1$ satisfy the monic recurrence $xh_n=h_{n+1}+\beta_n^2h_{n-1}$ with
$\beta_n^2=1/(4(n+\nu-1)(n+\nu))$, and their Jacobi operator has spectrum the closure of $\{\pm1/j_{\nu-1,k}\}$
\cite[\S3]{VA}; the orthogonality measure is symmetric with masses proportional to $j_{\nu-1,k}^{-2}$.
% TODO (citation): VA Sec. 3 is quoted only for the recurrence and the spectrum. The masses ~ j^{-2} follow from the
% partial fractions of the m-function -2 nu J_nu(1/x)/J_{nu-1}(1/x); this is done in PaperII Remark 9.6 only for
% nu = 3/2 (via DLMF 10.16.1, 4.22.3). For general nu quote the Mittag-Leffler expansion of J_nu/J_{nu-1}
% (DLMF 10.21-type) before submission; until then the identification is checked numerically (spec_lommel, 1e-15).
By the quadratic transformation of symmetric measures (as in \cite[Remark~9.6]{PaperII}), the image measure in
$y=x^2$ has $b_n^{(y)}=\beta_{2n-1}\beta_{2n}$, and $\sigma_\Lambda$ is its image under $y\mapsto\pi^2y$. Hence
$b_n=\pi^2/(4(2n+\mu)\sqrt{(2n+\mu)^2-1})$, and $c(2)=\pi^2/16$ gives the first identity. The second is
$(2N)\sqrt{4N^2-1}=4N^2(1-\frac1{4N^2})^{1/2}$.
\end{proof}

The first identity agrees with the Jacobi coefficients computed from the zeros (mpmath zeros for $k\le120$, McMahon's
expansion beyond, a continuous tail) to $1.6\cdot10^{-15}$ at $\mu=0$ and $\mu=1.5$, $n\le40$ (script
\texttt{spec\_lommel}). Here $\alpha=(s-1-p)/p=-\frac12$, and the large-$k$ form
$j_{\mu,k}=\pi(k+\frac\mu2-\frac14)+O(1/k)$ gives $Z_\Lambda(0)=\zeta(0,\frac\mu2+\frac34)=-\frac\mu2-\frac14$, so
$N=n+\frac\alpha2-Z_\Lambda(0)$: Proposition~\ref{prop:bessel} is the natural-index law of Section~\ref{sec:conj},
exactly.
% TODO (citation): McMahon's expansion (DLMF 10.21.19) is not quoted in the project's citation record.

\subsection{The residue recursion}\label{sec:recursion}
We use the frame of \S\ref{sec:frame} for $p\ne1$: the clockwise circles $\partial U_0$, $\partial U_r$ around $0$ and $4c$,
the local maps $f$ (hard edge) and $\tilde\varphi$ (soft edge), and the Szeg\H{o} quotients $\rho$, $\omega$ of
\cite[Sections~8--10]{PaperII}. Let $M:=\frac1{\sqrt2}\bigl(\begin{smallmatrix}1&i\\i&1\end{smallmatrix}\bigr)$, let $\circ$ be the
entrywise product, let $a_k=a_k(\alpha)$, $b_k=b_k(\alpha)$ be the coefficients of the Bessel expansions
\cite[(10.17.1), (10.17.8)]{DLMF} and $u_k,v_k$ those of the Airy expansions \cite[(9.7.2)]{DLMF}, and put
\[
  \Xi_k:=M^{-1}(M\circ C_k),\ C_k:=\begin{pmatrix}(-1)^ka_k&a_k\\(-1)^kb_k&b_k\end{pmatrix},\qquad
  X_k:=M^{-1}(M\circ C_k'),\ C_k':=\begin{pmatrix}u_k&(-1)^ku_k\\v_k&(-1)^kv_k\end{pmatrix}.
\]

\begin{lemma}[the jumps to every order]\label{lem:N6a}
Let $p\ne1$ and $K,J\ge1$. Uniformly on $\partial U_0\cup\partial U_r$, boundary values included, the jump of the error
matrix on the circles is $I+\sum_{k\le K}\Delta_kn^{-k}+O(n^{-K-1})$, where $\Delta_k$ equals
\begin{align*}
  \Delta_k^0&=(2f^{1/2})^{-k}P_0^\infty\rho^{\sigma_3}(\sigma_3\Xi_k\sigma_3)\rho^{-\sigma_3}(P_0^\infty)^{-1}&&\text{on }\partial U_0,\\
  \Delta_k^r&=(-1/\tilde\varphi)^kP_0^\infty(-\omega)^{-\sigma_3/2}X_k(-\omega)^{\sigma_3/2}(P_0^\infty)^{-1}&&\text{on }\partial U_r,
\end{align*}
and $\Delta_k=G_{\mu_n}\Delta_k^{\mathrm{mod}}G_{\mu_n}^{-1}+O(n^{-(J+1)p})$.
\end{lemma}

\begin{proof}[Proof sketch]
For $k=1$ this is \cite[Lemmas~9.1, 9.2 and~10.6]{PaperII}, and for $k=2$ the third-order version cited in
\S\ref{sec:E2E3}. The proof is the same with $K+1$ terms of the expansions \cite[(10.40.1)--(10.40.4), (10.17.5),
(10.17.6), (10.17.11), (10.17.12)]{DLMF} of the Bessel model \cite[\S6]{KMVV} and \cite[(9.7.5)--(9.7.6)]{DLMF} of the
Airy model, whose Poincar\'e remainders are $O(n^{-K-1})$ uniformly because the arguments are of size $n$ on the circles.
The sign pattern $(-1)^k$ is the same in every sector, for every $k$: each column of the Airy model is a constant
times $\mathrm{Ai}(Z)$ or $\mathrm{Ai}'(Z)$ with $Z=e^{2\pi im/3}\zeta$, $|\mathrm{ph}\,Z|\le\pi-\delta$, and
$Z^{3/2}=\pm\zeta^{3/2}$, so $(-1)^k$ sits on the column that carries $e^{-\xi}$ whichever branch occurs; the column
carrying $e^{\pm\xi}$ is fixed at $k=0$. The Bessel model is the same with $Y=\pm w$. The $n$-dependence of the weight is
Lemma~\ref{lem:Fprime} at order $J$, and the conjugation by $G_{\mu_n}$ is Lemma~\ref{lem:gauge}, whose proof commutes
diagonal factors and does not depend on $k$.
\end{proof}

\begin{lemma}[meromorphy]\label{lem:N6b}
For every $k$, $\Delta_k^{\mathrm{mod}}$ is meromorphic in $U_0$ and in $U_r$, and $\Delta_k^\mu:=G_\mu\Delta_k^{\mathrm{mod}}G_\mu^{-1}$
is meromorphic with Laurent coefficients that are entire functions of $\mu$.
\end{lemma}

\begin{proof}[Proof sketch]
For $p>1$: $\hat P^{\mathrm{mod}}(P_0^\infty)^{-1}$ is analytic in $U_e\setminus\{e\}$ and has the expansion of
Lemma~\ref{lem:N6a} uniformly on compact subsets, so $\Delta_k=\lim n^k(\cdots)$ is analytic and single-valued there; its
explicit formula grows at most polynomially at $e$, so $e$ is a pole. For $0<p<1$, where the glued hard-edge parametrix
carries seam jumps inside $U_0$, one checks directly that $\Delta_k$ has no jump: from the definitions,
$X_k=\frac12(u_k+v_k)I-\frac12(u_k-v_k)\sigma_2$ for even $k$ and $X_k=\frac12(u_k+v_k)\sigma_3+\frac12(v_k-u_k)i\sigma_1$ for odd
$k$, so $JX_kJ^{-1}=(-1)^kX_k$ for $J=i\sigma_2$; on $(4c-\epsilon,4c)$, $\tilde\varphi_+=-\tilde\varphi_-$ contributes the same
$(-1)^k$, and $P^\infty_{0,+}=P^\infty_{0,-}J$, $(-\omega)^{1/2}_+(-\omega)^{1/2}_-=1$. The hard edge is the same with
$\sigma_3\Xi_k\sigma_3$, $f^{1/2}_+=-f^{1/2}_-$ and the relations for $\rho$ of \cite[Section~10]{PaperII}. The gauge
$G_\mu$ is entire in $(z,\mu)$ and satisfies the same relations.
\end{proof}

For $\mu\in\C^J$ define, near each $e\in\{0,4c\}$, $\mathsf R_{\mathrm{in}}^{(0)}:=I$ and for $m\ge1$
\begin{equation}\label{eq:recursion}
\begin{gathered}
  \mathsf F_m(\mu):=\sum_{j=1}^m\mathsf R_{\mathrm{in}}^{(m-j)}(\mu)\Delta_j^\mu,\qquad
  \mathsf R_{\mathrm{out}}^{(m)}(\mu):=\sum_{e\in\{0,4c\}}\bigl(\text{principal part of }\mathsf F_m(\mu)\text{ at }e\bigr),\\
  \mathsf R_{\mathrm{in}}^{(m)}(\mu):=\mathsf R_{\mathrm{out}}^{(m)}(\mu)-\mathsf F_m(\mu),\qquad
  \mathsf R_1^{(m)}(\mu):=\Res_0\mathsf F_m(\mu)+\Res_{4c}\mathsf F_m(\mu).
\end{gathered}
\end{equation}
By Lemma~\ref{lem:N6b} these are well defined, $\mathsf R_{\mathrm{in}}^{(m)}$ is analytic in each disc, and all entries are
entire in $\mu$. For $m=1,2$ and $\mu=0$ they are the objects of \cite[Proposition~9.5]{PaperII} and its third-order
version.

\begin{proposition}[small norm to order $K$]\label{prop:N6c}
Let $p\ne1$ and $K,J\ge1$. Then $\mathsf R_1=\sum_{m\le K}\mathsf R_1^{(m)}(\mu_n)n^{-m}+O(n^{-K-1}+n^{-1-(J+1)p})$, and the
same expansion holds for $\mathsf R$ uniformly at positive distance from its jump contour.
\end{proposition}

\begin{proof}[Proof sketch]
Write $\eta_n:=n^{-K-1}+n^{-1-(J+1)p}$. By Lemma~\ref{lem:N6a} the jump is $I+\sum_{k\le K}\Delta_k^{\mu_n}n^{-k}$ on the circles
up to $O(\eta_n)$, and off the circles it is $I$ up to errors that are $O(\eta_n)$ in $L^1\cap L^2\cap L^\infty$, as in
\cite[Propositions~9.5 and~10.8]{PaperII}; the Cauchy operator is bounded uniformly in $n$. One proves by induction on
$m=1,\dots,K+1$ that the boundary value $\mathsf R_-$ equals $\sum_{k<m}\mathsf R_-^{(k)}n^{-k}+O_{L^2}(n^{-m}+n^{-1-(J+1)p})$, where
$\mathsf R^{(k)}$ is $\mathsf R^{(k)}_{\mathrm{in}}(\mu_n)$ inside the discs and $\mathsf R^{(k)}_{\mathrm{out}}(\mu_n)$ outside. In the
step, insert the claim into $\mathsf R=I+\mathcal C(\mathsf R_-(J-I))$: the new error is
$O(n^{-m-1}+n^{-1-(J+1)p})$, and the term of order $n^{-m}$ is the Cauchy integral over the clockwise circles of
$\mathsf F_m(\mu_n)$, which by the residue theorem is $\mathsf R_{\mathrm{out}}^{(m)}$ outside the discs and
$\mathsf R_{\mathrm{out}}^{(m)}-\mathsf F_m=\mathsf R_{\mathrm{in}}^{(m)}$ inside. The coefficient $\mathsf R_1$ is
$-\frac1{2\pi i}\int\mathsf R_-(J-I)$, and $-\frac1{2\pi i}\oint_{\mathrm{cw}}\mathsf F_m=\Res_0\mathsf F_m+\Res_{4c}\mathsf F_m$.
\end{proof}

Define $\cD_K(n;\mu):=((Q_1)_{12}(Q_1)_{21})^{1/2}/c$, expanded to order $n^{-K}$, where
$Q_1=P_1+\sum_{m\le K}\mathsf R_1^{(m)}(\mu)n^{-m}$, and write $\cD_K(n;\mu)=\sum_\beta\mu^\beta\cD_{\beta,K}(n)$ for its
Taylor series in $\mu$. Each $\cD_{\beta,K}$ is a polynomial in $1/n$ of degree at most $K$. A polynomial in $1/n$ is
\emph{written in natural form} when it is re-expanded in $1/N$ by $1/n=N^{-1}/(1-\delta_NN^{-1})$; a layer $\mu^\beta$ carries
$n^{-|\beta|_wp}=N^{-|\beta|_wp}(N/n)^{|\beta|_wp}$, with $|\beta|_w:=\sum_jj\beta_j$.

\begin{lemma}[exact algebra]\label{lem:N6d}
The recursion~\eqref{eq:recursion} and the extraction $\cD_K$ give, identically in the symbols:
\begin{enumerate}
\item[(a)] $\mu=0$, $K=6$: $\cD_{0,6}(n)=(n/N)^p[1+f_2N^{-2}+f_4N^{-4}+f_6N^{-6}+O(N^{-7})]$, with no $N^{-1},N^{-3},N^{-5}$
terms, $f_2$ as in~\eqref{eq:f2} and $f_4,f_6$ as in~\eqref{eq:f4}, \eqref{eq:f6}; in particular $d_1=-p(\alpha+1)/2$ and $d_2$ is
that of Theorem~\ref{thm:E2E3};
\item[(b)] first order in $\mu_j$, for every integer $j\ge1$ ($j$ a symbol), $K=4$: the coefficient of $n^{-1}$ per unit
$\binom{2j}j\mu_j$ is $pj/2$ (with $c=1$), and in natural form the layer is a multiple of
$N^{-1-jp}[1+g_jN^{-2}+O(N^{-4})]$, with no $N^{-1}$, $N^{-3}$ term in the bracket and $g_j=A_j+B_j\alpha^2$ as in
Theorem~\ref{thm:A}; the same holds at $K=5$ for the integers $j=1,\dots,5$, and for $j=1$ the $N^{-4}$ coefficient of
the bracket is $g_{1,4}$ of~\eqref{eq:g14};
\item[(c)] second order in $\mu_1$, $K=5$: per unit $\mu_1^2$ the natural form has no $N^{-1-2p}$ term, the coefficient
$p(3p+1)/2$ at $N^{-2-2p}$, no $N^{-3-2p}$ and no $N^{-5-2p}$ term, and $p^2h_2$ at $N^{-4-2p}$, with $h_2$ as in~\eqref{eq:h2}.
\end{enumerate}
\end{lemma}

\begin{proof}
With $c=1$ every input is a rational function of $(p,\alpha)$: the local maps have convergent series with rational
coefficients, $\pi$ cancels in $f$, the square roots are of series whose leading coefficient is $1$ or a perfect square,
and the factors $2^{-1/2}$ in $\mathsf T$ and $M$ occur in conjugate pairs. The recursion is carried out on truncated
Laurent series over $\mathbb Q(p,\alpha)[i]$; every series carries the index of its highest exactly known coefficient, and
every principal part checks that its coefficients down to the residue are known, otherwise the run stops. For (b),
$\mu_j$ is normalised by $\binom{2j}j$; the gauge to first order is $\partial_{\mu_j}\Delta_k^\mu|_{\mu=0}=[\mathcal G_j,\Delta_k^{\mathrm{mod}}]$
with $\mathcal G_j=\mathsf T\bigl(\begin{smallmatrix}0&-(z-4c)Q_j/2\\-zQ_j/2&0\end{smallmatrix}\bigr)\mathsf T^{-1}$, and it needs only the
Taylor data of $Q_j/\binom{2j}j$ at $0$ and at $4c=4$, which are rational in $j$: $\prod_{i\le m}(j-i)/(2(2j-2i-1))$ at $0$
and $j(j-1)\cdots(j-m)/(4^{m+1}(\frac12)_{m+1})$ at $4$, from $\sum_jQ_j(z)t^j=t/((1-zt)\sqrt{1-4t})$. So (b) is a
computation over $\mathbb Q(p,\alpha,j)[i]$. For (c) the gauge is taken to second order,
$\cosh h=1+z(z-4)\mu^2Q_1^2/8+O(\mu^4)$, $S=\mu Q_1/2+O(\mu^3)$. The scripts and their outputs are listed in
Appendix~\ref{app:repro}; each output begins with the SHA-256 hashes of the scripts that produced it. Since $(z-4c)Q_j$ and
$zQ_j$ scale as $c^j$ under $z=c\zeta$, $\mu_j$ enters only as $c^j\mu_j$; the integer layer and the ratios $g_j$, $h_2$
are independent of $c$.
\end{proof}

The closed forms are
\begin{equation}\label{eq:f4}
  f_4=\frac{(p+1)(6p^3+46p^2+55p+3)}{1440\,p(3p+2)}-\frac{5(p+1)(p+2)\alpha^2}{96(3p+2)}-\frac{p(p+2)^2\alpha^4}{128(3p+2)},
\end{equation}
\begin{equation}\label{eq:f6}
  f_6=\frac{\phi_0+\phi_2\alpha^2+\phi_4\alpha^4+\phi_6\alpha^6}{(3p+2)^2(5p+2)},
\end{equation}
\begin{align*}
  \phi_0&=-\frac{p+1}{362880\,p^2}\bigl(450p^8+1590p^7-316p^6-23761p^5-66864p^4-65507p^3\\
  &\qquad\qquad-18220p^2-1632p-40\bigr),\\
  \phi_2&=\frac{(p+1)(p+2)(90p^6+270p^5+664p^4-165p^3-1649p^2-288p+28)}{11520\,p},\\
  \phi_4&=-\frac{(p+1)(p+2)^2(20p^4+45p^3+162p^2+72p+16)}{1536},\\
  \phi_6&=\frac{p(p+2)^3(16p^3+7p^2-4p-4)}{3072},
\end{align*}
\begin{multline}\label{eq:g14}
  g_{1,4}=\frac{p+1}{5760\,p^2(3p+2)}\bigl(9p^6+117p^5+161p^4-569p^3-1550p^2-976p+48\bigr)\\
  +\frac{(p+1)(p+2)(3p^4+33p^3+124p^2+182p+108)\alpha^2}{192\,p(3p+2)}\\
  +\frac{(p+2)^3(p^3-4p-4)\alpha^4}{128\,p(3p+2)},
\end{multline}
\begin{equation}\label{eq:h2}
  h_2=\frac{(p+1)(p+2)(3p^2+2p+1)}{8p^2}+\frac{(p+2)(2p^3-p^2-7p-6)\alpha^2}{16p^2}.
\end{equation}
For $j=1$, $g_1=(p+1)(p^2+3p-2)/(24p)+(p+2)(p^2+2p+2)\alpha^2/(8p)$. At the Bessel point $(p,\alpha)=(2,-\frac12)$,
$f_4=\frac3{128}$ and $f_6=\frac5{1024}$, in agreement with Proposition~\ref{prop:bessel}.

\subsection{The theorem}\label{sec:N6}
\begin{theorem}\label{thm:N6}
Let $p>0$, $p\ne1$, $s>1$, and $K,J,B\ge1$.
\begin{enumerate}
\item[(i)] \emph{Double expansion.} $\displaystyle\frac{n^pb_n}c=\sum_{|\beta|\le B}\mu_n^\beta\cD_{\beta,K}(n)
  +O\bigl(n^{-K-1}+n^{-2-(B+1)p}+n^{-1-(J+1)p}\bigr)$.
\item[(ii)] \emph{Natural forms of the layers.} $\cD_{0,6}$, $\mu_{n,j}\cD_{e_j,4}$ and $\mu_{n,1}^2\cD_{2e_1,5}$ satisfy
\begin{gather*}
  \cD_{0,6}(n)=\Bigl(\frac nN\Bigr)^p\Bigl[1+\frac{f_2}{N^2}+\frac{f_4}{N^4}+\frac{f_6}{N^6}+O(N^{-7})\Bigr],\\
  \mu_{n,j}\cD_{e_j,4}(n)=\Bigl(\frac nN\Bigr)^p\frac{\varepsilon_j}{N^{1+jp}}\Bigl[1+\frac{g_j}{N^2}+O(N^{-4})\Bigr]\qquad(j\ge1),\\
  \mu_{n,1}^2\cD_{2e_1,5}(n)=\Bigl(\frac nN\Bigr)^p\frac{\varepsilon_1^2}{N^{2+2p}}\Bigl[\frac{3p+1}{2p}+\frac{h_2}{N^2}+O(N^{-4})\Bigr],
\end{gather*}
with the brackets free of odd powers of $1/N$ to the orders shown; for $j=1$ the second bracket is
$1+g_1N^{-2}+g_{1,4}N^{-4}+O(N^{-5})$.
\item[(iii)] \emph{The coefficients of $b_n$.} With $K=6$, $B=2$ and $J\ge6/p-1$, expansion~\eqref{eq:thmA} holds.
\end{enumerate}
\end{theorem}

\begin{proof}
(i) By~\eqref{eq:extract} and Proposition~\ref{prop:N6c}, $n^{2p}b_n^2=(Q_1)_{12}(Q_1)_{21}$ with
$Q_1=P_1+\sum_{m\le K}\mathsf R_1^{(m)}(\mu_n)n^{-m}+O(n^{-K-1}+n^{-1-(J+1)p})$. The $n^0$ term of $(Q_1)_{12}(Q_1)_{21}$ is
$c^2$, since $P_1$ does not depend on $\mu$, so each coefficient of $\cD_K(n;\mu)$ is a polynomial in the entries of the
$\mathsf R_1^{(i)}(\mu)$, $i\le K$, which are entire in $\mu$ and independent of $n$. Taylor's theorem in $\mu$ on $|\mu|\le1$
applies coefficientwise. The $n^0$ coefficient is $1$ for all $\mu$, and the $n^{-1}$ coefficient is $d_1+cpM_\mu(4c)$,
exactly linear in $\mu$ by Lemma~\ref{lem:E37}; hence $\cD_{\beta,K}=O(n^{-2})$ for $|\beta|\ge2$, and the Taylor remainder
after order $B$ is $O(|\mu|^{B+1}n^{-2})$. Since $|\mu_n|=O(n^{-p})$, this is $O(n^{-2-(B+1)p})$.

(ii) This is Lemma~\ref{lem:N6d} in natural form, with the normalisations
$\mu_{n,j}\binom{2j}jpc^jj/2=\varepsilon_jn^{-jp}$ and $\mu_{n,1}^2p(3p+1)/2\cdot c^2=\varepsilon_1^2n^{-2p}(3p+1)/(2p)$.

(iii) In~(i) keep $\beta=0$ at $K=6$; $\beta=e_j$ for every $j$ with $1+jp<\min(7,2+3p)$, at $K=5$ ($j\le5$) or $K=4$; and
$\beta=2e_1$ at $K=5$. The other terms with $|\beta|=2$ ($\mu_1\mu_2,\mu_2^2,\dots$) are $O(n^{-2-3p})$ because
$\cD_\beta=O(n^{-2})$ and $|\beta|_w\ge3$, and those with $|\beta|=3$ are $O(n^{-2-3p})$. The truncation errors are bounded as
follows. The $j=1$ bracket is known to relative order $N^{-5}$, an error $O(n^{-6-p})$, which is $O(n^{-2-3p})$ for $p<2$
and $O(n^{-8})$ for $p\ge2$. For $j\ge2$ the bracket error is $O(n^{-5-jp})$, and $5+2p>2+3p$ for $p<3$, $5+2p\ge11$ for
$p\ge3$. The $h_2$ term of the $\varepsilon_1^2$ layer is $O(n^{-4-2p})$, which is $O(n^{-2-3p})$ for $p<2$ and $O(n^{-7})$ for
$p\ge\frac32$. Finally $n^{-1-(J+1)p}=O(n^{-7})$.
\end{proof}

A term $CN^{-e}$ of~\eqref{eq:thmA} is a coefficient of the expansion of $n^pb_n/c$ (in powers of $N$, or equivalently,
after re-expansion, of $n$) exactly when $e<\min(7,2+3p)$. Table~\ref{tab:visible} lists the resulting ranges. The terms
outside their range are true identities of the layers in Theorem~\ref{thm:N6}(ii) but are not determined as coefficients
of $b_n$, because the uncomputed layers $\mu_1\mu_2$, $\mu_1^3$ enter at $n^{-2-3p}$.

\begin{table}[ht]
\centering
\small
\begin{tabular}{lll}
\toprule
term & exponent $e$ & a coefficient of $n^pb_n/c$ for\\
\midrule
$f_2N^{-2}$ & $2$ & all $p\ne1$\\
$f_4N^{-4}$ & $4$ & $p>\frac23$\\
$f_6N^{-6}$ & $6$ & $p>\frac43$\\
$\varepsilon_jN^{-1-jp}$ & $1+jp$ & $(j-3)p<1$ and $1+jp<7$\\
$\varepsilon_1g_1N^{-3-p}$ & $3+p$ & $\frac12<p<4$\\
$\varepsilon_2g_2N^{-3-2p}$ & $3+2p$ & $1<p<2$\\
$\varepsilon_1g_{1,4}N^{-5-p}$ & $5+p$ & $\frac32<p<2$\\
$\frac{3p+1}{2p}\varepsilon_1^2N^{-2-2p}$ & $2+2p$ & $p<\frac52$\\
$\varepsilon_jg_jN^{-3-jp}$ ($j\ge3$), $\varepsilon_1^2h_2N^{-4-2p}$ & & never at $B=2$\\
\bottomrule
\end{tabular}
\caption{Which terms of~\eqref{eq:thmA} are coefficients of $b_n$ ($p\ne1$). Coincident exponents merge: at $p=1/j_0$,
$\varepsilon_{j_0}N^{-2}$ merges with $f_2N^{-2}$, and $f_2+\varepsilon_{j_0}=f_2-p/12$, as in Theorem~\ref{thm:E2E3}.}
\label{tab:visible}
\end{table}

\begin{remark}\label{rem:N6}
(a) Theorem~\ref{thm:N6} does not prove evenness in $1/N$ to all orders, nor any closed form of the bare function; the
computed layers are listed in~(ii), and the layers $\mu_1\mu_2$, $\mu_2^2$, $\mu_1^3$ and higher are not computed. Nothing
is claimed at $p=1$, where the hard-edge and band-factor constructions of \cite{PaperII} change.
(b) The same recursion evaluated numerically at $40$ digits, with truncation depths $40$--$64$, gives $f_1=f_3=f_5=f_7=0$ to
$10^{-35}$ at every sampled $(p,\alpha)$ with $p\in[\frac13,12]$, each $f_{2k}$ even in $\alpha$, and
$f_8(3,0)=1311857/632311515$, $f_8(2,0)=3899/2^{21}$ after rational recognition. The denominators of $f_4$, $f_6$ are
$3p+2$ and $(3p+2)^2(5p+2)$; they come from the factors $1/(a'+k)$, $a'=\frac1p+\frac12$, of the hard-edge local map,
$a'+k=((2k+1)p+2)/(2p)$, and $632311515=3^5\cdot5\cdot11^3\cdot17\cdot23$ is $3^5\cdot5\,(3p+2)^3(5p+2)(7p+2)$ at $p=3$
(scripts \texttt{m4\_rh4}, \texttt{m4\_f8\_denominator}). These statements are numerical.
(c) Heuristically, evenness in $\alpha$ after the shift reflects the fact that the Bessel coefficients $a_k(\alpha)$,
$b_k(\alpha)$ depend on $\alpha$ only through $\alpha^2$, while elsewhere $\alpha$ enters only through the power of $z$ at the
hard edge, which the shift by $\frac\alpha2$ absorbs. We have no proof, and no explanation of the evenness in $1/N$.
(d) The factor $1/(2j-1)$ in $g_j$, read as a function of a continuous Mellin variable $w$, is a simple pole at
$w=\frac12$, the Mellin signature of $(1-u)^{+1/2}$, against $(1-u)^{-1/2}$ for the leading kernel
(Lemma~\ref{lem:K0}). Partial fractions give
$B_j=\frac{p+2}{8p}\bigl[\frac{p^2j}2+\frac{p(p+2)}4+\frac{(p+2)(p+4)}{4(2j-1)}\bigr]$. This suggests a second-order edge kernel
(Problem~\ref{prob:kernel2}).
\end{remark}
